% !TeX root = main.tex
\newpage
\section{Lighting Ball}

\subsection{Purpose}

\texttt{LightingBall} is the current lighting bring-up scenario after the
Cathedral and FreeStreamHetero studies were retired.  It keeps the useful
parts of the lighting work in a smaller, more inspectable scene: a stream of
mobile light particles approaches a spherical surface made visible by ordinary
boundary-marker particles.  The sphere gives a known three-dimensional shape
for testing lighting response, surface normals, reflection, and Vulkan export
parity.

\subsection{Configured Sphere}

The scenario uses the configuration item:
\begin{verbatim}
Lighting_ball = {
    x = x_center;
    y = y_center;
    z = z_center;
    radius = radius;
    material_id = material_id;
    wall_flag = wall_flag;
};
\end{verbatim}

The generator samples this sphere and inserts standard boundary particles at
the sampled locations.  These particles remain ordinary boundary markers:
they have boundary \texttt{ptype=2}, occupy cells for broad-phase
locality, and are not mobile physical objects.

The marker particles do not define the collision surface.  When a mobile
source discovers a local sphere marker, the dynamics evaluate the analytic
sphere equation from \texttt{Lighting\_ball}.  The sphere normal is:
\[
    \mathbf n =
    \frac{\mathbf x_{\mathrm{source}}-\mathbf c}
         {\lVert\mathbf x_{\mathrm{source}}-\mathbf c\rVert},
\]
where \(\mathbf c\) is the configured sphere center.  Physical penetration is
computed from the source radius and the signed distance to the analytic sphere
surface.  This keeps the boundary particles in the same role as wall markers:
they make nearby analytic geometry discoverable without becoming the geometry.

\subsection{Python Status}

The Python path is the working reference for the lighting ball.  The
generator \texttt{gbase/GenLightingBall.py} creates the sphere markers and
uses the shared streaming structure for the mobile particle stream.
\texttt{base/ForceDynamics.py} owns the analytic lighting-ball evaluator and
reflects a source from the configured sphere normal when the source is moving
into the surface.

Lighting display now uses the first boundary-space lighting rule.  Photons
remain transport/debug particles.  When a photon contacts the analytic sphere
or another boundary target, the contacted surface region stores the surface
material color in persistent boundary-space state.  Python remains the
behavioral reference for photon contact and material color selection; its
current display can use existing boundary markers as a coarse approximation.

The old \texttt{BoundaryLightSample}/\texttt{BoundaryLightState} current and
filtered path remains retired.  Eye position, field of view, diffuse terms,
and specular terms are not active lighting inputs.

\subsection{Vulkan Status}

\texttt{LightingBall\_Vk.inc} is aligned with the Python include at the
configuration level.  It selects the same sphere, cell domain, death bounds,
stream geometry, material properties, radius, and timestep, but writes Vulkan
output names such as \texttt{LightingBall\_Vk.tst}.

The current Vulkan simple export carries the lighting-ball sphere into the
runtime scene.  The generated test file writes \texttt{Lighting\_ball} as a
typed structure, and \texttt{ShaderObj} emits \texttt{sphere.glsl} with the
enabled state, center, radius, material id, and wall flag.  The simple exporter
emits \texttt{ForceDynamicsSimpleLightingBall.glsl}.

The basic particle applications use the clean \texttt{FPM.*} shader family.
\texttt{FPM.*} contains ordinary dynamics and particle diagnostics only.  It
does not include the lighting SSBO or photon lighting behavior.  The lighting
application uses the \texttt{FPML.*} shader family for photon-aware compute
and particle/debug rendering.  The lighting surface itself uses
\texttt{BoundaryLight.vert} and \texttt{BoundaryLight.frag}.  \texttt{FPML.comp}
deposits persistent surface color into the lighting resource, and the
lighting surface fragment shader reads that cell-owned state directly.

Vulkan rendering also uses the material-owned color-mode contract.  Material
records contain \texttt{color\_mode} and \texttt{color}.  The generated
\texttt{material.glsl} exposes \texttt{COLOR\_MODE\_COLLISION},
\texttt{COLOR\_MODE\_VELOCITY\_ANGLE}, \texttt{COLOR\_MODE\_SOLID}, and
\texttt{COLOR\_MODE\_LUMENS}.  The shared \texttt{color\_map.glsl} no longer
uses the old named solid color schemes.  In \texttt{LUMENS} mode, uncontacted
photons are transparent unless their photon material has
\texttt{debug\_visible=true}.  Surface illumination comes from boundary-space
state rather than from drawing the moving photon point as the surface color.
The retired eye, diffuse, and specular brightness controls are not part of
the active \texttt{LightingBall} contract.

Photon behavior is carried by runtime \texttt{ptype}, not by the current
\texttt{material\_id}.  During generation a stream whose base material has
\texttt{particle\_type="photon"} emits particles with photon \texttt{ptype}.
At the beginning of each compute frame the photon \texttt{material\_id} is
reset to the base photon material.  When the photon contacts the sphere or a
wall, the collision path sets the photon \texttt{material\_id} to the
contacted surface material for diagnostics and future boundary lighting.
The Python and Vulkan renderers treat \texttt{LUMENS} as debug photon
visibility, so this temporary material change does not become the final
visible surface lighting result.  No separate \texttt{hit\_material\_id}
particle field is part of the active ABI.

The active \texttt{LightingBall\_Py.inc} and
\texttt{LightingBall\_Vk.inc} files intentionally use only the ordinary
viewer controls:
\begin{verbatim}
view_center = [3.5, 2.5, 2.5];
zoom = 1.25;
gl_point_size = 2.0;
\end{verbatim}
The eye-aligned camera keys were removed from the active scene because the
eye no longer participates in lighting.  If that camera mode is needed later,
it should be reintroduced as view/camera configuration rather than as part of
the lighting model.

\subsection{Current Boundary-Space Contract}

\texttt{LightingBall} now proves the first boundary-space path in both Python
and Vulkan:
\begin{itemize}
    \item keep photons as transport/debug particles;
    \item map photon contacts to analytic surface hit cells;
    \item store persistent cell-owned color from the contacted surface
          material;
    \item render the analytic sphere surface from the owning cell color
          instead of drawing light only on the photon point;
    \item keep \texttt{color\_map.glsl} under exporter ownership so photon
          \texttt{ptype}, temporary photon \texttt{material\_id}, and material
          display behavior cannot drift from the generated Vulkan contract.
\end{itemize}

The current Vulkan lighting resource stores one \texttt{BoundaryLightRecord}
per spatial cell.  \texttt{ResourceLighting.cpp} sizes the buffer as
\texttt{CellAryW * CellAryH * CellAryL}.  The sphere render resource bakes the
owning cell id for each sphere vertex into the vertex data, and the lighting
fragment shader reads \texttt{BoundaryLight[surfaceCellId]}.  The validated
state is that \texttt{LightingBall} in Vulkan runs and shows solid cell-owned
surface lighting.  Spectral lighting is not yet working and is the next
lighting feature to revisit.
